\section{Semantic Fusion}

Semantic fusion provides a unified geometric representation for multiple
modalities such as text, image, and audio. Each modality is first encoded into
a vector space of dimension $d$, after which the fused representation is
constructed through a weighted geometric combination.

Let
\[
\mathbf{t},\ \mathbf{i},\ \mathbf{a} \in \mathbb{R}^d
\]
denote the text, image, and audio embeddings respectively. These embeddings
are produced by modality-specific encoders (e.g., transformer-based text
encoders, convolutional or vision-transformer image encoders, and spectral
audio encoders).

We define the fused semantic vector as:
\[
\mathbf{z} = W_t \mathbf{t} + W_i \mathbf{i} + W_a \mathbf{a}
\]
where
\[
W_t,\ W_i,\ W_a \in \mathbb{R}^{d \times d}
\]
are learnable geometric transformation matrices. These matrices align each
modality into a shared manifold, enabling cross-modal coherence.

To ensure geometric stability, the fused representation is normalized:
\[
\mathbf{z}_{norm} = \frac{\mathbf{z}}{\|\mathbf{z}\|}
\]

This normalization ensures that fused semantics remain bounded within the
manifold, preventing drift amplification and enabling consistent geodesic
evaluation across modalities.

\subsection{Geometric Interpretation}

The fusion process can be viewed as a projection into a shared semantic
manifold. Each modality contributes a directional component, and the weighted
sum defines a point whose curvature, drift, and gradient properties can be
evaluated by downstream geometric operators.

This geometric alignment is essential for the cognitive loop, as it allows
semantic routing, salience scoring, and stability filtering to operate on a
single unified representation.
